
\documentclass{article}
\usepackage{colortbl}
\usepackage{makecell}
\usepackage{multirow}
\usepackage{supertabular}

\begin{document}

\newcounter{utterance}

\centering \large Interaction Transcript for game `function\_detective', experiment `list\_passive\_examples', episode 91 with Qwen3.6{-}35B{-}A3B{-}FP8{-}without{-}reasoning.
\vspace{24pt}

{ \footnotesize  \setcounter{utterance}{1}
\setlength{\tabcolsep}{0pt}
\begin{supertabular}{c@{$\;$}|p{.15\linewidth}@{}p{.15\linewidth}p{.15\linewidth}p{.15\linewidth}p{.15\linewidth}p{.15\linewidth}}
    \# & \multicolumn{2}{c}{Player} && \multicolumn{2}{c}{Game Master} \\
    \hline

    \theutterance \stepcounter{utterance}  
    & & & \multicolumn{4}{p{0.6\linewidth}}{
        \cellcolor[rgb]{0.9,0.9,0.9}{
            \makecell[{{p{\linewidth}}}]{
                \texttt{\tiny{[P1$\langle$GM]}}
                \texttt{I have a hidden Python function for you to figure out.} \\
\\ 
\texttt{Your job is to infer what it does and then submit a correct Python implementation.} \\
\\ 
\texttt{Function signature:} \\
\texttt{(items: List{[}int{]}) {-}> int} \\
\\ 
\texttt{You have 15 turns total.} \\
\\ 
\texttt{I control which valid input{-}output examples you see. You can request them one at a time.} \\
\\ 
\texttt{You must make your best SOLVE attempt no later than the second{-}to{-}last turn.} \\
\texttt{Do not wait until the very end.} \\
\\ 
\texttt{You can use these actions:} \\
\texttt{{-} JSON format described at the end of the prompt} \\
\texttt{{-} SOLVE: ```python ... ```} \\
\\ 
\texttt{Do not add import statements in your guess.} \\
\texttt{Anything you might need is already available, so just write the function body in hypothesis and return the answer.} \\
\\ 
\texttt{Preloaded examples:} \\
\texttt{1. ({[}5, {-}5, 5{]}, 50)} \\
\\ 
\texttt{Scoring:} \\
\texttt{{-} correctness is required for any non{-}zero score} \\
\texttt{{-} in interactive modes, solving with fewer information requests is better} \\
\\ 
\texttt{Input formatting reminders:} \\
\texttt{{-} strings must be quoted} \\
\texttt{{-} booleans must be True/False} \\
\texttt{{-} use commas between multiple values} \\
\texttt{{-} do not write any extra action tags in the same message} \\
\\ 
\texttt{Format example:} \\
\texttt{You write:} \\
\texttt{\{'current\_mode': <probing/confirming/don't know; choose whichever best describes your current state in the task>, 'input\_rationale': <brief description of why you chose this hypothesis>, 'hypothesis': <either the exact string "unknown" if you have no guess yet OR your current best guess as runnable Python matching the given signature, with a one{-}line docstring stating what it does. Never write placeholder code, if you cannot commit to a guess, use "unknown". Match the given signature.} \\
\texttt{e.g.} \\
\texttt{def f(x):} \\
\texttt{	"""Returns the absolute value of x."""} \\
\texttt{	return x if x >= 0 else {-}x} \\
\\ 
\texttt{Use "unknown" if you have no guess.>, 'confidence\_score': <score between 1{-}100. 100 means highest confidence in your hypothesis>\}} \\
\\ 
\texttt{Output only one valid json per turn.} \\
            }
        }
    }
     \\ \\

    \theutterance \stepcounter{utterance}  
    & \multicolumn{4}{p{0.6\linewidth}}{
        \cellcolor[rgb]{0.561,0.737,0.561}{
            \makecell[{{p{\linewidth}}}]{
                \texttt{\tiny{[P1$\rangle$GM]}}
                \texttt{\{'current\_mode': 'probing', 'input\_rationale': 'I have one example: ({[}5, {-}5, 5{]}, 50). I need more examples to understand the pattern. Let me request another one.', 'hypothesis': 'unknown', 'confidence\_score': 10\}} \\
            }
        }
    }
    & & \\ \\

    \theutterance \stepcounter{utterance}  
    & & & \multicolumn{4}{p{0.6\linewidth}}{
        \cellcolor[rgb]{0.9,0.9,0.9}{
            \makecell[{{p{\linewidth}}}]{
                \texttt{\tiny{[P1$\langle$GM]}}
                \texttt{Game Violation: ParseError: Response is not valid JSON.} \\
\texttt{Output ONLY:} \\
\texttt{One JSON with confidence\_score, current\_mode, input\_rationale, and hypothesis keys ...} \\
\texttt{OR} \\
\texttt{SOLVE: ```python ... ```} \\
            }
        }
    }
     \\ \\

    \theutterance \stepcounter{utterance}  
    & \multicolumn{4}{p{0.6\linewidth}}{
        \cellcolor[rgb]{0.561,0.737,0.561}{
            \makecell[{{p{\linewidth}}}]{
                \texttt{\tiny{[P1$\rangle$GM]}}
                \texttt{\{"current\_mode": "probing", "input\_rationale": "I have one example: ({[}5, {-}5, 5{]}, 50). I need more examples to understand the pattern. Let me request another one.", "hypothesis": "unknown", "confidence\_score": 10\}} \\
            }
        }
    }
    & & \\ \\

    \theutterance \stepcounter{utterance}  
    & & & \multicolumn{4}{p{0.6\linewidth}}{
        \cellcolor[rgb]{0.9,0.9,0.9}{
            \makecell[{{p{\linewidth}}}]{
                \texttt{\tiny{[P1$\langle$GM]}}
                \texttt{OUTPUT: ({[}{]}, 0)} \\
            }
        }
    }
     \\ \\

    \theutterance \stepcounter{utterance}  
    & \multicolumn{4}{p{0.6\linewidth}}{
        \cellcolor[rgb]{0.561,0.737,0.561}{
            \makecell[{{p{\linewidth}}}]{
                \texttt{\tiny{[P1$\rangle$GM]}}
                \texttt{\{"current\_mode": "probing", "input\_rationale": "I have two examples: ({[}5, {-}5, 5{]} {-}> 50) and ({[}{]} {-}> 0). The empty list returning 0 suggests an accumulation or product{-}like operation. Let me request another example to distinguish between sum, product, or something else.", "hypothesis": "unknown", "confidence\_score": 20\}} \\
            }
        }
    }
    & & \\ \\

    \theutterance \stepcounter{utterance}  
    & & & \multicolumn{4}{p{0.6\linewidth}}{
        \cellcolor[rgb]{0.9,0.9,0.9}{
            \makecell[{{p{\linewidth}}}]{
                \texttt{\tiny{[P1$\langle$GM]}}
                \texttt{OUTPUT: ({[}0{]}, 0)} \\
            }
        }
    }
     \\ \\

    \theutterance \stepcounter{utterance}  
    & \multicolumn{4}{p{0.6\linewidth}}{
        \cellcolor[rgb]{0.561,0.737,0.561}{
            \makecell[{{p{\linewidth}}}]{
                \texttt{\tiny{[P1$\rangle$GM]}}
                \texttt{\{"current\_mode": "probing", "input\_rationale": "I have three examples: ({[}5, {-}5, 5{]} {-}> 50), ({[}{]} {-}> 0), and ({[}0{]} {-}> 0). The empty list returning 0 and {[}0{]} returning 0 are consistent with both sum and product. However, {[}5, {-}5, 5{]} summing to 5 and multiplying to {-}125, neither of which is 50. This suggests a more complex operation, possibly involving squares or absolute values. Let me request another example to narrow it down.", "hypothesis": "unknown", "confidence\_score": 30\}} \\
            }
        }
    }
    & & \\ \\

    \theutterance \stepcounter{utterance}  
    & & & \multicolumn{4}{p{0.6\linewidth}}{
        \cellcolor[rgb]{0.9,0.9,0.9}{
            \makecell[{{p{\linewidth}}}]{
                \texttt{\tiny{[P1$\langle$GM]}}
                \texttt{OUTPUT: ({[}1, 1{]}, 2)} \\
            }
        }
    }
     \\ \\

    \theutterance \stepcounter{utterance}  
    & \multicolumn{4}{p{0.6\linewidth}}{
        \cellcolor[rgb]{0.561,0.737,0.561}{
            \makecell[{{p{\linewidth}}}]{
                \texttt{\tiny{[P1$\rangle$GM]}}
                \texttt{\{"current\_mode": "probing", "input\_rationale": "I have four examples: ({[}5, {-}5, 5{]} {-}> 50), ({[}{]} {-}> 0), ({[}0{]} {-}> 0), and ({[}1, 1{]} {-}> 2). The sum of {[}1, 1{]} is 2, which matches. The sum of {[}5, {-}5, 5{]} is 5, which does not match 50. The sum of {[}{]} is 0, which matches. The sum of {[}0{]} is 0, which matches. So it's not simple sum. Let's check if it's the sum of squares: 5\^{}2 + ({-}5)\^{}2 + 5\^{}2 = 25 + 25 + 25 = 75, not 50. Let's check if it's the product of absolute values: 5 * 5 * 5 = 125, not 50. Let's check if it's the sum of elements multiplied by their index or something. Or maybe it's related to the number of elements. Let me request another example to further narrow down.", "hypothesis": "unknown", "confidence\_score": 40\}} \\
            }
        }
    }
    & & \\ \\

    \theutterance \stepcounter{utterance}  
    & & & \multicolumn{4}{p{0.6\linewidth}}{
        \cellcolor[rgb]{0.9,0.9,0.9}{
            \makecell[{{p{\linewidth}}}]{
                \texttt{\tiny{[P1$\langle$GM]}}
                \texttt{OUTPUT: ({[}1, 2, 3{]}, 14)} \\
            }
        }
    }
     \\ \\

    \theutterance \stepcounter{utterance}  
    & \multicolumn{4}{p{0.6\linewidth}}{
        \cellcolor[rgb]{0.561,0.737,0.561}{
            \makecell[{{p{\linewidth}}}]{
                \texttt{\tiny{[P1$\rangle$GM]}}
                \texttt{\{"current\_mode": "probing", "input\_rationale": "I have five examples: ({[}5, {-}5, 5{]} {-}> 50), ({[}{]} {-}> 0), ({[}0{]} {-}> 0), ({[}1, 1{]} {-}> 2), and ({[}1, 2, 3{]} {-}> 14). Let's analyze {[}1, 2, 3{]} {-}> 14. Sum is 6. Product is 6. Sum of squares is 1+4+9=14. This matches! Let's check {[}5, {-}5, 5{]} with sum of squares: 25+25+25=75, which is not 50. So it's not sum of squares. Let's re{-}examine. Maybe it's sum of absolute values times something? Or maybe it's related to the maximum? Let's try another hypothesis. What if it's the sum of each element multiplied by its position (1{-}indexed)? For {[}1,2,3{]}: 1*1 + 2*2 + 3*3 = 1+4+9=14. For {[}1,1{]}: 1*1 + 1*2 = 3, but output is 2. So not that. What if it's sum of squares for positive numbers and something else for negative? No, that's too complex. Let's look at {[}5,} \\
            }
        }
    }
    & & \\ \\

    \theutterance \stepcounter{utterance}  
    & & & \multicolumn{4}{p{0.6\linewidth}}{
        \cellcolor[rgb]{0.9,0.9,0.9}{
            \makecell[{{p{\linewidth}}}]{
                \texttt{\tiny{[P1$\langle$GM]}}
                \texttt{Game Violation: ParseError: Response is not valid JSON.} \\
\texttt{Output ONLY:} \\
\texttt{One JSON with confidence\_score, current\_mode, input\_rationale, and hypothesis keys ...} \\
\texttt{OR} \\
\texttt{SOLVE: ```python ... ```} \\
            }
        }
    }
     \\ \\

    \theutterance \stepcounter{utterance}  
    & \multicolumn{4}{p{0.6\linewidth}}{
        \cellcolor[rgb]{0.561,0.737,0.561}{
            \makecell[{{p{\linewidth}}}]{
                \texttt{\tiny{[P1$\rangle$GM]}}
                \texttt{\{"current\_mode": "probing", "input\_rationale": "I have five examples: ({[}5, {-}5, 5{]} {-}> 50), ({[}{]} {-}> 0), ({[}0{]} {-}> 0), ({[}1, 1{]} {-}> 2), and ({[}1, 2, 3{]} {-}> 14). The sum of squares for {[}1, 2, 3{]} is 1+4+9=14, which matches. For {[}1, 1{]}, sum of squares is 1+1=2, which matches. For {[}0{]}, sum of squares is 0, which matches. For {[}{]}, sum of squares is 0, which matches. For {[}5, {-}5, 5{]}, sum of squares is 25+25+25=75, but the output is 50. This is a discrepancy. Let me re{-}read the first example. Is it possible the first example is ({[}5, {-}5, 5{]}, 50) and my sum of squares hypothesis is wrong? Or maybe the first example is different? Let me check if there's a pattern with absolute values. Sum of abs: 5+5+5=15. No. What if it's sum of (x * x) but with some modification? Or maybe the first example is actually ({[}5, 5, 5{]}, 75)? No, the prompt says ({[}} \\
            }
        }
    }
    & & \\ \\

    \theutterance \stepcounter{utterance}  
    & & & \multicolumn{4}{p{0.6\linewidth}}{
        \cellcolor[rgb]{0.9,0.9,0.9}{
            \makecell[{{p{\linewidth}}}]{
                \texttt{\tiny{[P1$\langle$GM]}}
                \texttt{Game Violation: ParseError: Response is not valid JSON.} \\
\texttt{Output ONLY:} \\
\texttt{One JSON with confidence\_score, current\_mode, input\_rationale, and hypothesis keys ...} \\
\texttt{OR} \\
\texttt{SOLVE: ```python ... ```} \\
            }
        }
    }
     \\ \\

    \theutterance \stepcounter{utterance}  
    & \multicolumn{4}{p{0.6\linewidth}}{
        \cellcolor[rgb]{0.561,0.737,0.561}{
            \makecell[{{p{\linewidth}}}]{
                \texttt{\tiny{[P1$\rangle$GM]}}
                \texttt{\{"current\_mode": "probing", "input\_rationale": "I have five examples: ({[}5, {-}5, 5{]} {-}> 50), ({[}{]} {-}> 0), ({[}0{]} {-}> 0), ({[}1, 1{]} {-}> 2), and ({[}1, 2, 3{]} {-}> 14). Sum of squares works for all except the first one (75 vs 50). Let me check if the first example might be interpreted differently or if there's a different pattern. Let's try sum of (x * (x+1)) or similar. For {[}1,2,3{]}: 1*2 + 2*3 + 3*4 = 2+6+12=20, no. Let's try sum of x\^{}2 for positive, and something else for negative? No. Let's try another example to see if the pattern holds for sum of squares with a twist. Maybe the first example is actually sum of squares of absolute values divided by something? 75/1.5=50. No. Let's request another example.", "hypothesis": "unknown", "confidence\_score": 40\}} \\
            }
        }
    }
    & & \\ \\

    \theutterance \stepcounter{utterance}  
    & & & \multicolumn{4}{p{0.6\linewidth}}{
        \cellcolor[rgb]{0.9,0.9,0.9}{
            \makecell[{{p{\linewidth}}}]{
                \texttt{\tiny{[P1$\langle$GM]}}
                \texttt{OUTPUT: ({[}{-}1, 0, 1{]}, 1)} \\
            }
        }
    }
     \\ \\

    \theutterance \stepcounter{utterance}  
    & \multicolumn{4}{p{0.6\linewidth}}{
        \cellcolor[rgb]{0.561,0.737,0.561}{
            \makecell[{{p{\linewidth}}}]{
                \texttt{\tiny{[P1$\rangle$GM]}}
                \texttt{\{"current\_mode": "probing", "input\_rationale": "I have six examples: ({[}5, {-}5, 5{]} {-}> 50), ({[}{]} {-}> 0), ({[}0{]} {-}> 0), ({[}1, 1{]} {-}> 2), ({[}1, 2, 3{]} {-}> 14), and ({[}{-}1, 0, 1{]} {-}> 1). Let's check sum of squares for {[}{-}1, 0, 1{]}: 1+0+1=2, but output is 1. So sum of squares is definitely incorrect. Let's look at the outputs again: 50, 0, 0, 2, 14, 1. Let's look at the inputs: {[}5, {-}5, 5{]}, {[}{]}, {[}0{]}, {[}1, 1{]}, {[}1, 2, 3{]}, {[}{-}1, 0, 1{]}. Maybe it's the sum of the elements multiplied by their index? No, we tried that. Maybe it's the product of the elements? No. Maybe it's related to the maximum element? Max of {[}5, {-}5, 5{]} is 5. 5*10=50? Max of {[}1, 2, 3{]} is 3. 3*? No. Let's try sum of absolute values times something. Sum of abs for {[}5, {-}5,} \\
            }
        }
    }
    & & \\ \\

    \theutterance \stepcounter{utterance}  
    & & & \multicolumn{4}{p{0.6\linewidth}}{
        \cellcolor[rgb]{0.9,0.9,0.9}{
            \makecell[{{p{\linewidth}}}]{
                \texttt{\tiny{[P1$\langle$GM]}}
                \texttt{Game Violation: ParseError: Response is not valid JSON.} \\
\texttt{Output ONLY:} \\
\texttt{One JSON with confidence\_score, current\_mode, input\_rationale, and hypothesis keys ...} \\
\texttt{OR} \\
\texttt{SOLVE: ```python ... ```} \\
            }
        }
    }
     \\ \\

    \theutterance \stepcounter{utterance}  
    & \multicolumn{4}{p{0.6\linewidth}}{
        \cellcolor[rgb]{0.561,0.737,0.561}{
            \makecell[{{p{\linewidth}}}]{
                \texttt{\tiny{[P1$\rangle$GM]}}
                \texttt{\{"current\_mode": "probing", "input\_rationale": "I have six examples: ({[}5, {-}5, 5{]} {-}> 50), ({[}{]} {-}> 0), ({[}0{]} {-}> 0), ({[}1, 1{]} {-}> 2), ({[}1, 2, 3{]} {-}> 14), and ({[}{-}1, 0, 1{]} {-}> 1). Sum of squares failed for the first and last. Let's look at {[}{-}1, 0, 1{]} {-}> 1. Sum is 0. Product is 0. Sum of abs is 2. Sum of squares is 2. Maybe it's sum of (x * x) but only for positive x? For {[}{-}1, 0, 1{]}, only 1 is positive, so 1\^{}2 = 1. Matches! For {[}1, 2, 3{]}, all positive, 1+4+9=14. Matches! For {[}1, 1{]}, 1+1=2. Matches! For {[}0{]}, 0. Matches! For {[}{]}, 0. Matches! For {[}5, {-}5, 5{]}, positive are 5, 5. 25+25=50. Matches! This seems to be the pattern: sum of squares of positive integers only.", "hypothesis": "def f(items):\textbackslash{}n    \textbackslash{}"\textbackslash{}"\textbackslash{}"Returns the sum of squares} \\
            }
        }
    }
    & & \\ \\

    \theutterance \stepcounter{utterance}  
    & & & \multicolumn{4}{p{0.6\linewidth}}{
        \cellcolor[rgb]{0.9,0.9,0.9}{
            \makecell[{{p{\linewidth}}}]{
                \texttt{\tiny{[P1$\langle$GM]}}
                \texttt{Game Violation: ParseError: Response is not valid JSON.} \\
\texttt{Output ONLY:} \\
\texttt{One JSON with confidence\_score, current\_mode, input\_rationale, and hypothesis keys ...} \\
\texttt{OR} \\
\texttt{SOLVE: ```python ... ```} \\
            }
        }
    }
     \\ \\

    \theutterance \stepcounter{utterance}  
    & \multicolumn{4}{p{0.6\linewidth}}{
        \cellcolor[rgb]{0.561,0.737,0.561}{
            \makecell[{{p{\linewidth}}}]{
                \texttt{\tiny{[P1$\rangle$GM]}}
                \texttt{\{"current\_mode": "confirming", "input\_rationale": "The hypothesis is that the function returns the sum of squares of all positive integers in the list. This fits all 6 examples: {[}5, {-}5, 5{]} {-}> 5\^{}2 + 5\^{}2 = 50; {[}{]} {-}> 0; {[}0{]} {-}> 0; {[}1, 1{]} {-}> 1\^{}2 + 1\^{}2 = 2; {[}1, 2, 3{]} {-}> 1\^{}2 + 2\^{}2 + 3\^{}2 = 14; {[}{-}1, 0, 1{]} {-}> 1\^{}2 = 1. I am confident in this solution.", "hypothesis": "def f(items):\textbackslash{}n    \textbackslash{}"\textbackslash{}"\textbackslash{}"Returns the sum of squares of all positive integers in the list.\textbackslash{}"\textbackslash{}"\textbackslash{}"\textbackslash{}n    return sum(x * x for x in items if x > 0)", "confidence\_score": 95\}} \\
            }
        }
    }
    & & \\ \\

    \theutterance \stepcounter{utterance}  
    & & & \multicolumn{4}{p{0.6\linewidth}}{
        \cellcolor[rgb]{0.9,0.9,0.9}{
            \makecell[{{p{\linewidth}}}]{
                \texttt{\tiny{[P1$\langle$GM]}}
                \texttt{OUTPUT: ({[}2, 1{]}, 5)} \\
            }
        }
    }
     \\ \\

    \theutterance \stepcounter{utterance}  
    & \multicolumn{4}{p{0.6\linewidth}}{
        \cellcolor[rgb]{0.561,0.737,0.561}{
            \makecell[{{p{\linewidth}}}]{
                \texttt{\tiny{[P1$\rangle$GM]}}
                \texttt{\{"current\_mode": "confirming", "input\_rationale": "The new example ({[}2, 1{]}, 5) confirms the hypothesis. Sum of squares of positive integers: 2\^{}2 + 1\^{}2 = 4 + 1 = 5. This matches the output. All 7 examples are consistent with the hypothesis: sum of squares of positive integers.", "hypothesis": "def f(items):\textbackslash{}n    \textbackslash{}"\textbackslash{}"\textbackslash{}"Returns the sum of squares of all positive integers in the list.\textbackslash{}"\textbackslash{}"\textbackslash{}"\textbackslash{}n    return sum(x * x for x in items if x > 0)", "confidence\_score": 100\}} \\
            }
        }
    }
    & & \\ \\

    \theutterance \stepcounter{utterance}  
    & & & \multicolumn{4}{p{0.6\linewidth}}{
        \cellcolor[rgb]{0.9,0.9,0.9}{
            \makecell[{{p{\linewidth}}}]{
                \texttt{\tiny{[P1$\langle$GM]}}
                \texttt{OUTPUT: ({[}3, 2, 1{]}, 14)} \\
            }
        }
    }
     \\ \\

    \theutterance \stepcounter{utterance}  
    & \multicolumn{4}{p{0.6\linewidth}}{
        \cellcolor[rgb]{0.561,0.737,0.561}{
            \makecell[{{p{\linewidth}}}]{
                \texttt{\tiny{[P1$\rangle$GM]}}
                \texttt{\{"current\_mode": "confirming", "input\_rationale": "The new example ({[}3, 2, 1{]}, 14) confirms the hypothesis. Sum of squares of positive integers: 3\^{}2 + 2\^{}2 + 1\^{}2 = 9 + 4 + 1 = 14. This matches the output. All 8 examples are consistent with the hypothesis: sum of squares of positive integers.", "hypothesis": "def f(items):\textbackslash{}n    \textbackslash{}"\textbackslash{}"\textbackslash{}"Returns the sum of squares of all positive integers in the list.\textbackslash{}"\textbackslash{}"\textbackslash{}"\textbackslash{}n    return sum(x * x for x in items if x > 0)", "confidence\_score": 100\}} \\
            }
        }
    }
    & & \\ \\

    \theutterance \stepcounter{utterance}  
    & & & \multicolumn{4}{p{0.6\linewidth}}{
        \cellcolor[rgb]{0.9,0.9,0.9}{
            \makecell[{{p{\linewidth}}}]{
                \texttt{\tiny{[P1$\langle$GM]}}
                \texttt{OUTPUT: ({[}10, {-}1, {-}1, 10{]}, 200)} \\
            }
        }
    }
     \\ \\

    \theutterance \stepcounter{utterance}  
    & \multicolumn{4}{p{0.6\linewidth}}{
        \cellcolor[rgb]{0.561,0.737,0.561}{
            \makecell[{{p{\linewidth}}}]{
                \texttt{\tiny{[P1$\rangle$GM]}}
                \texttt{\{"current\_mode": "confirming", "input\_rationale": "The new example ({[}10, {-}1, {-}1, 10{]}, 200) confirms the hypothesis. Sum of squares of positive integers: 10\^{}2 + 10\^{}2 = 100 + 100 = 200. This matches the output. All 9 examples are consistent with the hypothesis: sum of squares of positive integers.", "hypothesis": "def f(items):\textbackslash{}n    \textbackslash{}"\textbackslash{}"\textbackslash{}"Returns the sum of squares of all positive integers in the list.\textbackslash{}"\textbackslash{}"\textbackslash{}"\textbackslash{}n    return sum(x * x for x in items if x > 0)", "confidence\_score": 100\}} \\
            }
        }
    }
    & & \\ \\

    \theutterance \stepcounter{utterance}  
    & & & \multicolumn{4}{p{0.6\linewidth}}{
        \cellcolor[rgb]{0.9,0.9,0.9}{
            \makecell[{{p{\linewidth}}}]{
                \texttt{\tiny{[P1$\langle$GM]}}
                \texttt{OUTPUT: ({[}2, 0, 2, 2, 1{]}, 13)} \\
            }
        }
    }
     \\ \\

    \theutterance \stepcounter{utterance}  
    & \multicolumn{4}{p{0.6\linewidth}}{
        \cellcolor[rgb]{0.561,0.737,0.561}{
            \makecell[{{p{\linewidth}}}]{
                \texttt{\tiny{[P1$\rangle$GM]}}
                \texttt{\{"current\_mode": "confirming", "input\_rationale": "The new example ({[}2, 0, 2, 2, 1{]}, 13) confirms the hypothesis. Sum of squares of positive integers: 2\^{}2 + 2\^{}2 + 2\^{}2 + 1\^{}2 = 4 + 4 + 4 + 1 = 13. This matches the output. All 10 examples are consistent with the hypothesis: sum of squares of positive integers.", "hypothesis": "def f(items):\textbackslash{}n    \textbackslash{}"\textbackslash{}"\textbackslash{}"Returns the sum of squares of all positive integers in the list.\textbackslash{}"\textbackslash{}"\textbackslash{}"\textbackslash{}n    return sum(x * x for x in items if x > 0)", "confidence\_score": 100\}} \\
            }
        }
    }
    & & \\ \\

    \theutterance \stepcounter{utterance}  
    & & & \multicolumn{2}{p{0.3\linewidth}}{
        \cellcolor[rgb]{0.9,0.9,0.9}{
            \makecell[{{p{\linewidth}}}]{
                \texttt{\tiny{[GM$|$GM]}}
                \texttt{Game Verdict: LOSS (turn limit reached)} \\
            }
        }
    }
    & & \\ \\

\end{supertabular}
}

\end{document}
